% ===== RESUMEN =====
\clearpage
\chapter*{Resumen}
\addcontentsline{toc}{chapter}{Resumen}
\markboth{Resumen}{Resumen}

\noindent\textbf{TÍTULO:} ESTUDIO DE LAS ABERRACIONES VECTORIALES EN SISTEMAS
ESTIGMÁTICOS\footnote{Trabajo de grado.}

\vspace{0.5em}
\noindent\textbf{AUTOR:} JHONATAN S. BLANCO\footnote{Facultad de Ciencias.
Escuela de Física. Director: Rafael Ángel Torres Amarís, Físico, Mg., Dr.
Codirector: Cristian Hernández Celis, Físico, Mg.}

\vspace{0.5em}
\noindent\textbf{PALABRAS CLAVE:} difracción vectorial, aberraciones
vectoriales, superficies cartesianas, sistemas estigmáticos, coeficientes de
Fresnel, polarización, función de dispersión puntual.

\vspace{0.5em}
\noindent\textbf{DESCRIPCIÓN:}

En este trabajo se calculó el campo eléctrico en el plano focal de un sistema
estigmático simple, formado por una única superficie cartesiana que separa
dos medios homogéneos, mediante un formalismo vectorial que da cuenta de la
geometría de la superficie a través de su radio de curvatura y de la
propiedad estigmática, e incorpora su transmisión de Fresnel en una función
pupila vectorial generalizada, reducida a transformadas de Hankel bajo
iluminación azimutalmente simétrica. Se evidenció así la existencia de
aberraciones de origen vectorial en los sistemas estigmáticos: aun cuando la
aberración geométrica es nula, las condiciones de frontera de Maxwell~\cite{jackson1999} inducen
en la refracción una transformación anisótropa del campo, distinta en cada
punto de la superficie. Sobre el plano transversal la superficie actúa como
un diatenuador cuyos modos propios son las polarizaciones radial y azimutal,
modulados por los coeficientes de Fresnel $t_p$ y $t_s$; la aberración
vectorial, definida como la desviación respecto del límite escalar
$t_p=t_s=1$ en el que se recupera la teoría escalar de la difracción, aparece
modulada por el factor $t_-=t_p-t_s$, no nulo para cualquier interfaz entre
medios de distinto índice. Como manifestaciones se obtuvieron la deformación
del punto focal, caracterizada por dos radios principales que siguen a la
polarización lineal incidente; la aparición de polarización cruzada con
patrón de cuatro lóbulos; el enfoque de la componente longitudinal; y la
formación anisótropa de imágenes, ilustrada con un ejemplo litográfico. El
formalismo se extiende de manera natural a la reflexión, a sistemas de varias
superficies por composición de operadores y a iluminaciones arbitrarias.

% ===== ABSTRACT =====
\clearpage
\chapter*{Abstract}
\addcontentsline{toc}{chapter}{Abstract}
\markboth{Abstract}{Abstract}

\noindent\textbf{TITLE:} STUDY OF VECTORIAL ABERRATIONS IN STIGMATIC
SYSTEMS\footnote{Bachelor's thesis.}

\vspace{0.5em}
\noindent\textbf{AUTHOR:} JHONATAN S. BLANCO\footnote{Facultad de Ciencias.
Escuela de Física. Director: Rafael Ángel Torres Amarís, Físico, Mg., Dr.
Codirector: Cristian Hernández Celis, Físico, Mg.}

\vspace{0.5em}
\noindent\textbf{KEYWORDS:} vector diffraction, vectorial aberrations,
Cartesian surfaces, stigmatic systems, Fresnel coefficients, polarization,
point spread function.

\vspace{0.5em}
\noindent\textbf{DESCRIPTION:}

We compute the focal electric field of stigmatic refracting surfaces and show
they act as diattenuators with radial and azimuthal eigenmodes, producing
focal deformation, cross polarization, and longitudinal focusing despite
vanishing geometrical aberration. The field at the focal plane of a single
Cartesian surface separating two homogeneous media is obtained through a
vectorial formalism that accounts for the geometry of the surface through its
radius of curvature and the stigmatic property, and carries its Fresnel
transmission into a generalized vectorial pupil function, reduced to Hankel
transforms under azimuthally symmetric illumination. This evidences
the existence of aberrations of purely vectorial origin in stigmatic systems:
even though the stigmatic condition removes the optical-path aberration,
Maxwell's boundary conditions~\cite{jackson1999} induce an anisotropic, position-dependent
transformation of the field upon refraction. The vectorial aberration,
defined as the deviation from the scalar limit $t_p=t_s=1$ in which scalar
diffraction theory is exactly recovered, is modulated by the factor
$t_-=t_p-t_s$, nonvanishing for any interface between media of different
refractive index. Its manifestations include the deformation of the focal
spot, characterized by two principal radii that follow the incident linear
polarization; the appearance of a four-lobed cross-polarized component; the
focusing of the longitudinal component; and anisotropic image formation,
illustrated with a lithographic example. The formalism extends naturally to
reflection, to multi-surface systems by operator composition, and to
arbitrary illuminations.
